\documentclass[11pt]{article}
\usepackage[margin=1in]{geometry}
\usepackage{amsmath, amssymb, amsthm}
\usepackage{hyperref}
\usepackage{enumitem}

\title{Section 4.1: Vector Spaces and Subspaces \\ Comprehensive Problem Bank}
\author{}
\date{}

\begin{document}

\maketitle
\tableofcontents
\newpage

% ---------------------------------------------------------
\section{Concept Checklist \& Question Mapping}
% ---------------------------------------------------------
This section covers the foundational concepts of abstract vector spaces and subspaces, mirroring the archetypes, tricks, and pitfalls identified in the source material.

\begin{itemize}
    \item \textbf{Concept 1: Counterexample Generation (Non-Subspaces)}
    \begin{itemize}
        \item Identifying failure of the Zero Vector test.
        \item Identifying failure of Closure under Addition (often due to non-linear conditions like $xy=0$ or unions).
        \item Identifying failure of Closure under Scalar Multiplication (often due to inequalities like $x \le 0$).
    \end{itemize}
    \item \textbf{Concept 2: Subspace Verification via Axioms (The 3-Step Test)}
    \begin{itemize}
        \item Applying the rigorous 3-step test to subsets of $\mathbb{R}^n$.
        \item Applying the test to abstract spaces: Matrices ($M_{m \times n}$), Polynomials ($\mathbb{P}_n$), and Functions ($C[a,b]$).
    \end{itemize}
    \item \textbf{Concept 3: Span Identification (Parameter Factoring)}
    \begin{itemize}
        \item Translating a set defined by parameterized variables into a formal Span of constant vectors/matrices/polynomials.
    \end{itemize}
    \item \textbf{Concept 4: Span Membership via Linear Systems}
    \begin{itemize}
        \item Setting up augmented matrices to determine if a vector lies in a given span.
        \item Using row reduction (echelon form) to check for system consistency.
    \end{itemize}
    \item \textbf{Concept 5: Theoretical Properties and Synthesis}
    \begin{itemize}
        \item Proving properties about intersections and sums of subspaces.
        \item Abstract proofs relying on the 10 vector space axioms.
    \end{itemize}
    \item \textbf{Concept 6: Diagnostic Testing ("Find the Flaw")}
    \begin{itemize}
        \item Identifying logical fallacies (e.g., proving by specific example).
        \item Identifying algebraic misinterpretations (e.g., confusing scalar $0$ with the zero vector, or checking matrix multiplication instead of addition).
    \end{itemize}
\end{itemize}

\vspace{0.5cm}
\noindent \textbf{Question Mapping:}
\begin{itemize}
    \item \textbf{Concept 1:} Problems 1, 2, 3, 4, 5, 6
    \item \textbf{Concept 2:} Problems 7, 8, 9, 10, 11, 12
    \item \textbf{Concept 3:} Problems 13, 14, 15, 16, 17
    \item \textbf{Concept 4:} Problems 18, 19, 20, 21, 22
    \item \textbf{Concept 5:} Problems 23, 24, 25, 26
    \item \textbf{Concept 6:} Problems 27, 28, 29, 30
\end{itemize}

\newpage

% ---------------------------------------------------------
\section{Problem Sets}
% ---------------------------------------------------------

\subsection*{Concept 1: Counterexample Generation (Non-Subspaces)}
\textit{Determine if the following sets are subspaces of their respective vector spaces. If not, provide a specific counterexample showing which of the 3 subspace rules fails.}

\begin{enumerate}
    \item Let $V = \left\{ \begin{bmatrix} x \\ y \end{bmatrix} \in \mathbb{R}^2 : x \le 0, y \le 0 \right\}$. Is $V$ a subspace of $\mathbb{R}^2$?
    \item Let $W = \left\{ \begin{bmatrix} x \\ y \end{bmatrix} \in \mathbb{R}^2 : x^2 + y^2 \le 1 \right\}$. Is $W$ a subspace of $\mathbb{R}^2$?
    \item Let $H$ be the set of all polynomials in $\mathbb{P}_2$ such that $p(1) = 2$. Is $H$ a subspace of $\mathbb{P}_2$?
    \item Let $K = \{A \in M_{2\times 2} : \det(A) = 0\}$. Is $K$ a subspace of $M_{2\times 2}$?
    \item Let $U = \left\{ \begin{bmatrix} x \\ y \\ z \end{bmatrix} \in \mathbb{R}^3 : xyz = 0 \right\}$. Is $U$ a subspace of $\mathbb{R}^3$?
    \item Let $S$ be the set of all continuous functions $f \in C[0,1]$ such that $f(x) \ge 0$ for all $x \in [0,1]$. Is $S$ a subspace of $C[0,1]$?
\end{enumerate}

\subsection*{Concept 2: Subspace Verification via Axioms (The 3-Step Test)}
\textit{Use the rigorous 3-step subspace test to prove that the following sets are subspaces of their respective vector spaces.}

\begin{enumerate}[resume]
    \item Let $H = \left\{ \begin{bmatrix} x \\ y \\ z \end{bmatrix} \in \mathbb{R}^3 : 2x - 3y + z = 0 \right\}$. Prove $H$ is a subspace of $\mathbb{R}^3$.
    \item Let $W = \{A \in M_{2\times 2} : A^T = A\}$. (The set of all symmetric $2 \times 2$ matrices). Prove $W$ is a subspace of $M_{2\times 2}$.
    \item Let $V = \{p \in \mathbb{P}_3 : p(0) = 0 \text{ and } p(2) = 0\}$. Prove $V$ is a subspace of $\mathbb{P}_3$.
    \item Let $H = \{f \in C(\mathbb{R}) : f(-x) = -f(x) \text{ for all } x\}$. (The set of all odd functions). Prove $H$ is a subspace of $C(\mathbb{R})$.
    \item Let $W = \left\{ \begin{bmatrix} a & b \\ c & d \end{bmatrix} \in M_{2\times 2} : a+d = 0 \right\}$. Prove $W$ is a subspace of $M_{2\times 2}$.
    \item Let $U = \left\{ \begin{bmatrix} x \\ y \\ z \\ w \end{bmatrix} \in \mathbb{R}^4 : x=y \text{ and } z=w \right\}$. Prove $U$ is a subspace of $\mathbb{R}^4$.
\end{enumerate}

\subsection*{Concept 3: Span Identification (Parameter Factoring)}
\textit{For each of the following sets, factor out the parameters to express the set as the Span of a finite set of vectors/matrices/polynomials. Conclude that it is a subspace via Theorem 1.}

\begin{enumerate}[resume]
    \item $H = \left\{ \begin{bmatrix} 2s - t \\ 3s \\ t \end{bmatrix} : s, t \in \mathbb{R} \right\}$.
    \item $W = \{ a + bt + (a-b)t^2 : a, b \in \mathbb{R} \}$ in $\mathbb{P}_2$.
    \item $V = \left\{ \begin{bmatrix} x \\ y \\ z \end{bmatrix} \in \mathbb{R}^3 : x - 2y + 5z = 0 \right\}$. (Hint: Solve for one variable in terms of the others first).
    \item $H = \left\{ \begin{bmatrix} a+2b & a \\ b & a-b \end{bmatrix} : a, b \in \mathbb{R} \right\}$ in $M_{2\times 2}$.
    \item $U = \{ p \in \mathbb{P}_2 : p(t) = ct^2, c \in \mathbb{R} \}$.
\end{enumerate}

\subsection*{Concept 4: Span Membership via Linear Systems}
\textit{Use augmented matrices and row reduction to answer the following questions.}

\begin{enumerate}[resume]
    \item Is $y = \begin{bmatrix} 1 \\ 2 \\ 3 \end{bmatrix}$ in the span of $v_1 = \begin{bmatrix} 1 \\ 1 \\ 0 \end{bmatrix}, v_2 = \begin{bmatrix} 0 \\ 1 \\ 1 \end{bmatrix}, v_3 = \begin{bmatrix} 1 \\ 0 \\ -1 \end{bmatrix}$?
    \item For what value(s) of $h$ is $y = \begin{bmatrix} 2 \\ -1 \\ h \end{bmatrix}$ in the span of $v_1 = \begin{bmatrix} 1 \\ 0 \\ 2 \end{bmatrix}$ and $v_2 = \begin{bmatrix} -1 \\ 1 \\ -1 \end{bmatrix}$?
    \item Is the polynomial $p(t) = 1 - t + 2t^2$ in the span of $p_1(t) = 1+t$, $p_2(t) = 1-t^2$, and $p_3(t) = t+t^2$?
    \item Is the matrix $A = \begin{bmatrix} 1 & 2 \\ 2 & 1 \end{bmatrix}$ in the span of $M_1 = \begin{bmatrix} 1 & 0 \\ 0 & 1 \end{bmatrix}$ and $M_2 = \begin{bmatrix} 0 & 1 \\ 1 & 0 \end{bmatrix}$?
    \item Determine if the vector $y = \begin{bmatrix} 4 \\ -2 \\ 1 \end{bmatrix}$ is in the subspace spanned by $v_1 = \begin{bmatrix} 1 \\ -1 \\ 1 \end{bmatrix}$ and $v_2 = \begin{bmatrix} 2 \\ 0 \\ -1 \end{bmatrix}$.
\end{enumerate}

\subsection*{Concept 5: Theoretical Properties and Synthesis}
\textit{Provide a formal proof for each of the following statements.}

\begin{enumerate}[resume]
    \item Prove that if $H, K,$ and $L$ are all subspaces of a vector space $V$, then their intersection $H \cap K \cap L$ is also a subspace of $V$.
    \item Let $H$ be a subspace of $V$. Prove that if $u \in H$ and $u+v \in H$, then $v$ must also be in $H$.
    \item Let $v_1, v_2$ be vectors in a vector space $V$. Use the 3-step test to prove that $H = \text{Span}\{v_1, v_2\}$ is a subspace of $V$.
    \item Let $H$ and $K$ be subspaces of $V$. Define the sum of the subspaces as $H+K = \{h+k : h \in H, k \in K\}$. Prove that $H+K$ is a subspace of $V$.
\end{enumerate}

\subsection*{Concept 6: Diagnostic Testing ("Find the Flaw")}
\textit{Each of the following solutions contains a single, subtle conceptual or algebraic flaw. Identify the error, explain why it is wrong, and provide the correct solution.}

\begin{enumerate}[resume]
    \item \textbf{Question:} Let $H = \left\{ \begin{bmatrix} x \\ y \end{bmatrix} : x+y = 0 \right\}$. Is $H$ a subspace of $\mathbb{R}^2$? \\
    \textbf{Flawed Solution:} We check closure under addition. Let $u = \begin{bmatrix} 1 \\ -1 \end{bmatrix}$ and $v = \begin{bmatrix} -2 \\ 2 \end{bmatrix}$. Both are in $H$ because $1+(-1)=0$ and $-2+2=0$. Their sum is $u+v = \begin{bmatrix} -1 \\ 1 \end{bmatrix}$. Since $-1+1=0$, the sum is in $H$. Therefore, $H$ is closed under addition.
    
    \item \textbf{Question:} Let $H$ be the set of all invertible $2 \times 2$ matrices. Is $H$ a subspace of $M_{2\times 2}$? \\
    \textbf{Flawed Solution:} Let $A$ and $B$ be in $H$. Since they are invertible, $\det(A) \neq 0$ and $\det(B) \neq 0$. We check closure: $\det(AB) = \det(A)\det(B)$. Since neither determinant is zero, their product is not zero. Thus $AB$ is invertible and is in $H$. $H$ is a subspace.
    
    \item \textbf{Question:} Let $H = \{ p \in \mathbb{P}_2 : p(2) = 0 \}$. Is $H$ a subspace? \\
    \textbf{Flawed Solution:} To be a subspace, $H$ must contain the zero vector. The zero vector is the number $0$. Since $0$ is a number and not a polynomial, $H$ does not contain the zero vector and is not a subspace.
    
    \item \textbf{Question:} Let $H = \left\{ \begin{bmatrix} x \\ y \end{bmatrix} : x^2 = y^2 \right\}$. Is $H$ a subspace of $\mathbb{R}^2$? \\
    \textbf{Flawed Solution:} We check scalar multiplication. Let $c$ be a scalar and $u = \begin{bmatrix} x \\ y \end{bmatrix} \in H$, meaning $x^2 = y^2$. For $cu = \begin{bmatrix} cx \\ cy \end{bmatrix}$, we check the condition: $(cx)^2 = c^2x^2$. Since $x^2=y^2$, this equals $c^2y^2 = (cy)^2$. Thus, $cu \in H$. Because it is closed under scalar multiplication, it is a subspace.
\end{enumerate}

\newpage

% ---------------------------------------------------------
\section{Detailed Solutions}
% ---------------------------------------------------------

\subsection*{Concept 1: Counterexample Generation}

\textbf{Solution 1:}
No, $V$ is not a subspace. While it contains the zero vector and is closed under addition, it fails closure under scalar multiplication. 
\textit{Counterexample:} Let $u = \begin{bmatrix} -1 \\ -1 \end{bmatrix} \in V$ (since $-1 \le 0$). Let $c = -1$. Then $cu = \begin{bmatrix} 1 \\ 1 \end{bmatrix}$. Since $1 \not\le 0$, $cu \notin V$.

\vspace{0.3cm}
\textbf{Solution 2:}
No, $W$ is not a subspace. The condition $x^2+y^2 \le 1$ represents the unit disk, which is bounded. It fails closure under scalar multiplication.
\textit{Counterexample:} Let $u = \begin{bmatrix} 1 \\ 0 \end{bmatrix} \in W$ (since $1^2+0^2 = 1 \le 1$). Let $c = 2$. Then $cu = \begin{bmatrix} 2 \\ 0 \end{bmatrix}$. Checking the condition: $2^2 + 0^2 = 4 \not\le 1$. Thus $cu \notin W$.

\vspace{0.3cm}
\textbf{Solution 3:}
No, $H$ is not a subspace. We use the "Zero Vector Cheat". The zero vector in $\mathbb{P}_2$ is the zero polynomial $0(t) = 0$ for all $t$. Evaluating this at $t=1$ gives $0(1) = 0$. Since $0 \neq 2$, the zero polynomial is not in $H$.

\vspace{0.3cm}
\textbf{Solution 4:}
No, $K$ is not a subspace. It fails closure under addition.
\textit{Counterexample:} Let $A = \begin{bmatrix} 1 & 0 \\ 0 & 0 \end{bmatrix}$ and $B = \begin{bmatrix} 0 & 0 \\ 0 & 1 \end{bmatrix}$. Both are in $K$ because $\det(A) = 0$ and $\det(B) = 0$. However, $A+B = \begin{bmatrix} 1 & 0 \\ 0 & 1 \end{bmatrix} = I$. The determinant of $I$ is $1 \neq 0$, so $A+B \notin K$.

\vspace{0.3cm}
\textbf{Solution 5:}
No, $U$ is not a subspace. The condition $xyz=0$ means at least one coordinate must be zero (the union of the three coordinate planes). It fails closure under addition.
\textit{Counterexample:} Let $u = \begin{bmatrix} 1 \\ 1 \\ 0 \end{bmatrix} \in U$ (since $1\cdot 1\cdot 0 = 0$) and $v = \begin{bmatrix} 0 \\ 0 \\ 1 \end{bmatrix} \in U$ (since $0\cdot 0\cdot 1 = 0$). Their sum is $u+v = \begin{bmatrix} 1 \\ 1 \\ 1 \end{bmatrix}$. The product of the components is $1\cdot 1\cdot 1 = 1 \neq 0$. Thus $u+v \notin U$.

\vspace{0.3cm}
\textbf{Solution 6:}
No, $S$ is not a subspace. It fails closure under scalar multiplication.
\textit{Counterexample:} Let $f(x) = 1$ for all $x$. Since $1 \ge 0$, $f \in S$. Let $c = -1$. The new function is $(cf)(x) = -1$. Since $-1 \not\ge 0$, $cf \notin S$.

\subsection*{Concept 2: Subspace Verification via Axioms}

\textbf{Solution 7:}
We apply the 3-step test.
1. \textbf{Zero Vector:} Let $x=0, y=0, z=0$. $2(0) - 3(0) + 0 = 0$. Thus $\begin{bmatrix} 0 \\ 0 \\ 0 \end{bmatrix} \in H$.
2. \textbf{Addition:} Let $u = \begin{bmatrix} x_1 \\ y_1 \\ z_1 \end{bmatrix}, v = \begin{bmatrix} x_2 \\ y_2 \\ z_2 \end{bmatrix} \in H$. This means $2x_1 - 3y_1 + z_1 = 0$ and $2x_2 - 3y_2 + z_2 = 0$.
$u+v = \begin{bmatrix} x_1+x_2 \\ y_1+y_2 \\ z_1+z_2 \end{bmatrix}$. Checking the condition: $2(x_1+x_2) - 3(y_1+y_2) + (z_1+z_2) = (2x_1 - 3y_1 + z_1) + (2x_2 - 3y_2 + z_2) = 0 + 0 = 0$. Thus $u+v \in H$.
3. \textbf{Scalar Multiplication:} Let $c \in \mathbb{R}$. $cu = \begin{bmatrix} cx_1 \\ cy_1 \\ cz_1 \end{bmatrix}$. Checking the condition: $2(cx_1) - 3(cy_1) + (cz_1) = c(2x_1 - 3y_1 + z_1) = c(0) = 0$. Thus $cu \in H$. $H$ is a subspace.

\vspace{0.3cm}
\textbf{Solution 8:}
1. \textbf{Zero Vector:} The $2 \times 2$ zero matrix $0$ satisfies $0^T = 0$, so $0 \in W$.
2. \textbf{Addition:} Let $A, B \in W$, so $A^T = A$ and $B^T = B$. We check $A+B$: $(A+B)^T = A^T + B^T = A + B$. Thus $A+B \in W$.
3. \textbf{Scalar Multiplication:} Let $c \in \mathbb{R}$. We check $cA$: $(cA)^T = c(A^T) = cA$. Thus $cA \in W$. $W$ is a subspace.

\vspace{0.3cm}
\textbf{Solution 9:}
1. \textbf{Zero Vector:} The zero polynomial $0(t) = 0$ satisfies $0(0) = 0$ and $0(2) = 0$. Thus $0 \in V$.
2. \textbf{Addition:} Let $p, q \in V$. Then $(p+q)(0) = p(0) + q(0) = 0 + 0 = 0$. Also $(p+q)(2) = p(2) + q(2) = 0 + 0 = 0$. Thus $p+q \in V$.
3. \textbf{Scalar Multiplication:} Let $c \in \mathbb{R}$. $(cp)(0) = c \cdot p(0) = c(0) = 0$. Also $(cp)(2) = c \cdot p(2) = c(0) = 0$. Thus $cp \in V$. $V$ is a subspace.

\vspace{0.3cm}
\textbf{Solution 10:}
1. \textbf{Zero Vector:} The zero function $z(x) = 0$ satisfies $z(-x) = 0 = -0 = -z(x)$. Thus $z \in H$.
2. \textbf{Addition:} Let $f, g \in H$. $(f+g)(-x) = f(-x) + g(-x) = -f(x) - g(x) = -(f(x) + g(x)) = -(f+g)(x)$. Thus $f+g \in H$.
3. \textbf{Scalar Multiplication:} Let $c \in \mathbb{R}$. $(cf)(-x) = c \cdot f(-x) = c(-f(x)) = -(cf)(x)$. Thus $cf \in H$. $H$ is a subspace.

\vspace{0.3cm}
\textbf{Solution 11:}
1. \textbf{Zero Vector:} For the zero matrix, $a=0, d=0$. $0+0=0$, so $0 \in W$.
2. \textbf{Addition:} Let $A = \begin{bmatrix} a_1 & b_1 \\ c_1 & d_1 \end{bmatrix}, B = \begin{bmatrix} a_2 & b_2 \\ c_2 & d_2 \end{bmatrix} \in W$. Then $a_1+d_1=0$ and $a_2+d_2=0$. $A+B = \begin{bmatrix} a_1+a_2 & b_1+b_2 \\ c_1+c_2 & d_1+d_2 \end{bmatrix}$. The sum of the diagonal entries is $(a_1+a_2) + (d_1+d_2) = (a_1+d_1) + (a_2+d_2) = 0 + 0 = 0$. Thus $A+B \in W$.
3. \textbf{Scalar Multiplication:} Let $k \in \mathbb{R}$. $kA = \begin{bmatrix} ka_1 & kb_1 \\ kc_1 & kd_1 \end{bmatrix}$. The sum of the diagonal entries is $ka_1 + kd_1 = k(a_1+d_1) = k(0) = 0$. Thus $kA \in W$. $W$ is a subspace.

\vspace{0.3cm}
\textbf{Solution 12:}
1. \textbf{Zero Vector:} $\begin{bmatrix} 0 \\ 0 \\ 0 \\ 0 \end{bmatrix}$ satisfies $0=0$ and $0=0$.
2. \textbf{Addition:} Let $u = (x_1, x_1, z_1, z_1)$ and $v = (x_2, x_2, z_2, z_2)$. $u+v = (x_1+x_2, x_1+x_2, z_1+z_2, z_1+z_2)$. The first two entries are equal, and the last two are equal. Thus $u+v \in U$.
3. \textbf{Scalar Multiplication:} $cu = (cx_1, cx_1, cz_1, cz_1)$. The first two entries are equal, and the last two are equal. Thus $cu \in U$. $U$ is a subspace.

\subsection*{Concept 3: Span Identification}

\textbf{Solution 13:}
We factor the parameters $s$ and $t$ out of the vector:
$\begin{bmatrix} 2s - t \\ 3s \\ t \end{bmatrix} = \begin{bmatrix} 2s \\ 3s \\ 0 \end{bmatrix} + \begin{bmatrix} -t \\ 0 \\ t \end{bmatrix} = s\begin{bmatrix} 2 \\ 3 \\ 0 \end{bmatrix} + t\begin{bmatrix} -1 \\ 0 \\ 1 \end{bmatrix}$.
Thus, $H = \text{Span}\left\{ \begin{bmatrix} 2 \\ 3 \\ 0 \end{bmatrix}, \begin{bmatrix} -1 \\ 0 \\ 1 \end{bmatrix} \right\}$. By Theorem 1, it is a subspace.

\vspace{0.3cm}
\textbf{Solution 14:}
We group the terms by the parameters $a$ and $b$:
$p(t) = a + bt + at^2 - bt^2 = a(1 + t^2) + b(t - t^2)$.
Thus, $W = \text{Span}\{ 1+t^2, t-t^2 \}$. By Theorem 1, it is a subspace.

\vspace{0.3cm}
\textbf{Solution 15:}
The condition is $x - 2y + 5z = 0$. We can solve for $x$: $x = 2y - 5z$.
Substitute this back into a generic vector in $V$:
$\begin{bmatrix} x \\ y \\ z \end{bmatrix} = \begin{bmatrix} 2y - 5z \\ y \\ z \end{bmatrix} = y\begin{bmatrix} 2 \\ 1 \\ 0 \end{bmatrix} + z\begin{bmatrix} -5 \\ 0 \\ 1 \end{bmatrix}$.
Thus, $V = \text{Span}\left\{ \begin{bmatrix} 2 \\ 1 \\ 0 \end{bmatrix}, \begin{bmatrix} -5 \\ 0 \\ 1 \end{bmatrix} \right\}$. By Theorem 1, it is a subspace.

\vspace{0.3cm}
\textbf{Solution 16:}
We factor the parameters $a$ and $b$ out of the matrix:
$\begin{bmatrix} a+2b & a \\ b & a-b \end{bmatrix} = a\begin{bmatrix} 1 & 1 \\ 0 & 1 \end{bmatrix} + b\begin{bmatrix} 2 & 0 \\ 1 & -1 \end{bmatrix}$.
Thus, $H = \text{Span}\left\{ \begin{bmatrix} 1 & 1 \\ 0 & 1 \end{bmatrix}, \begin{bmatrix} 2 & 0 \\ 1 & -1 \end{bmatrix} \right\}$. By Theorem 1, it is a subspace.

\vspace{0.3cm}
\textbf{Solution 17:}
The polynomial is already factored by the parameter $c$:
$p(t) = c(t^2)$.
Thus, $U = \text{Span}\{ t^2 \}$. By Theorem 1, it is a subspace.

\subsection*{Concept 4: Span Membership}

\textbf{Solution 18:}
We set up the augmented matrix $[v_1 \ v_2 \ v_3 \ | \ y]$:
$\begin{bmatrix} 1 & 0 & 1 & | & 1 \\ 1 & 1 & 0 & | & 2 \\ 0 & 1 & -1 & | & 3 \end{bmatrix}$.
Row reduce: $R_2 \leftarrow R_2 - R_1$:
$\begin{bmatrix} 1 & 0 & 1 & | & 1 \\ 0 & 1 & -1 & | & 1 \\ 0 & 1 & -1 & | & 3 \end{bmatrix}$.
$R_3 \leftarrow R_3 - R_2$:
$\begin{bmatrix} 1 & 0 & 1 & | & 1 \\ 0 & 1 & -1 & | & 1 \\ 0 & 0 & 0 & | & 2 \end{bmatrix}$.
The last row is $[0 \ 0 \ 0 \ | \ 2]$, which means the system is inconsistent. Therefore, $y$ is \textbf{not} in the span.

\vspace{0.3cm}
\textbf{Solution 19:}
Set up the augmented matrix:
$\begin{bmatrix} 1 & -1 & | & 2 \\ 0 & 1 & | & -1 \\ 2 & -1 & | & h \end{bmatrix}$.
$R_3 \leftarrow R_3 - 2R_1$:
$\begin{bmatrix} 1 & -1 & | & 2 \\ 0 & 1 & | & -1 \\ 0 & 1 & | & h-4 \end{bmatrix}$.
$R_3 \leftarrow R_3 - R_2$:
$\begin{bmatrix} 1 & -1 & | & 2 \\ 0 & 1 & | & -1 \\ 0 & 0 & | & h-3 \end{bmatrix}$.
For the system to be consistent, we must have $h-3 = 0 \implies h = 3$.

\vspace{0.3cm}
\textbf{Solution 20:}
We equate coefficients: $c_1(1+t) + c_2(1-t^2) + c_3(t+t^2) = 1 - t + 2t^2$.
Constants: $c_1 + c_2 = 1$
$t$ terms: $c_1 + c_3 = -1$
$t^2$ terms: $-c_2 + c_3 = 2$
Augmented matrix:
$\begin{bmatrix} 1 & 1 & 0 & | & 1 \\ 1 & 0 & 1 & | & -1 \\ 0 & -1 & 1 & | & 2 \end{bmatrix}$.
$R_2 \leftarrow R_2 - R_1$:
$\begin{bmatrix} 1 & 1 & 0 & | & 1 \\ 0 & -1 & 1 & | & -2 \\ 0 & -1 & 1 & | & 2 \end{bmatrix}$.
$R_3 \leftarrow R_3 - R_2$:
$\begin{bmatrix} 1 & 1 & 0 & | & 1 \\ 0 & -1 & 1 & | & -2 \\ 0 & 0 & 0 & | & 4 \end{bmatrix}$.
The system is inconsistent. The polynomial is \textbf{not} in the span.

\vspace{0.3cm}
\textbf{Solution 21:}
We want $c_1 M_1 + c_2 M_2 = A$.
$c_1 \begin{bmatrix} 1 & 0 \\ 0 & 1 \end{bmatrix} + c_2 \begin{bmatrix} 0 & 1 \\ 1 & 0 \end{bmatrix} = \begin{bmatrix} c_1 & c_2 \\ c_2 & c_1 \end{bmatrix}$.
We need this to equal $\begin{bmatrix} 1 & 2 \\ 2 & 1 \end{bmatrix}$.
By matching entries, $c_1 = 1$ and $c_2 = 2$. Since a solution exists, $A$ \textbf{is} in the span.

\vspace{0.3cm}
\textbf{Solution 22:}
Augmented matrix:
$\begin{bmatrix} 1 & 2 & | & 4 \\ -1 & 0 & | & -2 \\ 1 & -1 & | & 1 \end{bmatrix}$.
$R_2 \leftarrow R_2 + R_1$, $R_3 \leftarrow R_3 - R_1$:
$\begin{bmatrix} 1 & 2 & | & 4 \\ 0 & 2 & | & 2 \\ 0 & -3 & | & -3 \end{bmatrix}$.
$R_2 \leftarrow \frac{1}{2}R_2$:
$\begin{bmatrix} 1 & 2 & | & 4 \\ 0 & 1 & | & 1 \\ 0 & -3 & | & -3 \end{bmatrix}$.
$R_3 \leftarrow R_3 + 3R_2$:
$\begin{bmatrix} 1 & 2 & | & 4 \\ 0 & 1 & | & 1 \\ 0 & 0 & | & 0 \end{bmatrix}$.
The system is consistent (solution is $c_2=1, c_1=2$). Yes, $y$ is in the subspace.

\subsection*{Concept 5: Theoretical Properties}

\textbf{Solution 23:}
1. \textbf{Zero Vector:} Since $H, K, L$ are subspaces, $0 \in H$, $0 \in K$, and $0 \in L$. Therefore, $0 \in H \cap K \cap L$.
2. \textbf{Addition:} Let $u, v \in H \cap K \cap L$. This means $u,v \in H$, $u,v \in K$, and $u,v \in L$. Because they are subspaces, $u+v \in H$, $u+v \in K$, and $u+v \in L$. Thus $u+v \in H \cap K \cap L$.
3. \textbf{Scalar Multiplication:} Let $u \in H \cap K \cap L$ and $c \in \mathbb{R}$. Then $u \in H, K, L$. Because they are subspaces, $cu \in H, K, L$. Thus $cu \in H \cap K \cap L$. It is a subspace.

\vspace{0.3cm}
\textbf{Solution 24:}
Since $H$ is a subspace, it is closed under scalar multiplication. Because $u \in H$, the scalar multiple $(-1)u = -u$ is also in $H$.
Since $H$ is closed under addition, the sum of any two vectors in $H$ is in $H$. We are given $u+v \in H$, and we know $-u \in H$.
Adding them together: $(u+v) + (-u) = v$. Therefore, $v \in H$.

\vspace{0.3cm}
\textbf{Solution 25:}
$H = \{ c_1 v_1 + c_2 v_2 : c_1, c_2 \in \mathbb{R} \}$.
1. \textbf{Zero Vector:} Let $c_1=0, c_2=0$. Then $0v_1 + 0v_2 = 0 \in H$.
2. \textbf{Addition:} Let $u = a_1 v_1 + a_2 v_2$ and $w = b_1 v_1 + b_2 v_2$ be in $H$.
$u+w = (a_1 v_1 + a_2 v_2) + (b_1 v_1 + b_2 v_2) = (a_1+b_1)v_1 + (a_2+b_2)v_2$. Since $a_1+b_1$ and $a_2+b_2$ are real scalars, $u+w \in H$.
3. \textbf{Scalar Multiplication:} Let $k \in \mathbb{R}$. $ku = k(a_1 v_1 + a_2 v_2) = (ka_1)v_1 + (ka_2)v_2$. Since $ka_1$ and $ka_2$ are real scalars, $ku \in H$. $H$ is a subspace.

\vspace{0.3cm}
\textbf{Solution 26:}
1. \textbf{Zero Vector:} Since $H$ and $K$ are subspaces, $0 \in H$ and $0 \in K$. Thus $0+0 = 0 \in H+K$.
2. \textbf{Addition:} Let $u, v \in H+K$. Then $u = h_1+k_1$ and $v = h_2+k_2$ for some $h_1, h_2 \in H$ and $k_1, k_2 \in K$.
$u+v = (h_1+k_1) + (h_2+k_2) = (h_1+h_2) + (k_1+k_2)$. Since $H$ and $K$ are closed under addition, $h_1+h_2 \in H$ and $k_1+k_2 \in K$. Thus $u+v \in H+K$.
3. \textbf{Scalar Multiplication:} Let $c \in \mathbb{R}$. $cu = c(h_1+k_1) = ch_1 + ck_1$. Since $H$ and $K$ are closed under scalar multiplication, $ch_1 \in H$ and $ck_1 \in K$. Thus $cu \in H+K$. $H+K$ is a subspace.

\subsection*{Concept 6: Diagnostic Testing ("Find the Flaw")}

\textbf{Solution 27:}
\textbf{Flaw:} The student attempted to prove closure under addition by picking specific numerical vectors. Proving a set IS a subspace requires showing it holds for \textit{all} generic vectors in the set.
\textbf{Correct Solution:} Let $u = \begin{bmatrix} x_1 \\ y_1 \end{bmatrix}$ and $v = \begin{bmatrix} x_2 \\ y_2 \end{bmatrix}$ be in $H$. This means $x_1+y_1=0$ and $x_2+y_2=0$.
$u+v = \begin{bmatrix} x_1+x_2 \\ y_1+y_2 \end{bmatrix}$. We check the condition: $(x_1+x_2) + (y_1+y_2) = (x_1+y_1) + (x_2+y_2) = 0 + 0 = 0$. Thus $u+v \in H$. (Note: $H$ is indeed a subspace, but the proof method was flawed).

\vspace{0.3cm}
\textbf{Solution 28:}
\textbf{Flaw:} The student checked closure under matrix \textit{multiplication} ($AB$). Vector spaces require closure under matrix \textit{addition} ($A+B$).
\textbf{Correct Solution:} $H$ is not a subspace. It fails the zero vector test because the zero matrix is not invertible ($\det(0) = 0$). It also fails closure under addition: Let $A = I$ and $B = -I$. Both are invertible, but $A+B = 0$, which is not invertible.

\vspace{0.3cm}
\textbf{Solution 29:}
\textbf{Flaw:} The student confused the scalar number $0$ with the zero vector in $\mathbb{P}_2$, which is the zero polynomial $p(t) = 0$.
\textbf{Correct Solution:} The zero vector in $\mathbb{P}_2$ is $p(t) = 0$ for all $t$. We check if it satisfies the condition $p(2) = 0$. Yes, evaluating the zero polynomial at $t=2$ gives $0$. Thus, $H$ \textit{does} contain the zero vector. (It is, in fact, a subspace).

\vspace{0.3cm}
\textbf{Solution 30:}
\textbf{Flaw:} The student only checked scalar multiplication and concluded it was a subspace, completely skipping the check for closure under addition.
\textbf{Correct Solution:} $H$ is not a subspace because it fails closure under addition. Let $u = \begin{bmatrix} 1 \\ 1 \end{bmatrix} \in H$ (since $1^2=1^2$) and $v = \begin{bmatrix} 1 \\ -1 \end{bmatrix} \in H$ (since $1^2=(-1)^2$). Their sum is $u+v = \begin{bmatrix} 2 \\ 0 \end{bmatrix}$. Checking the condition: $2^2 = 4$ and $0^2 = 0$. Since $4 \neq 0$, $u+v \notin H$.

\newpage

% ---------------------------------------------------------
\section*{Concept Checklist \& Question Mapping (Reference)}
% ---------------------------------------------------------
\begin{itemize}
    \item \textbf{Concept 1: Counterexample Generation (Non-Subspaces)} $\rightarrow$ Problems 1, 2, 3, 4, 5, 6
    \item \textbf{Concept 2: Subspace Verification via Axioms (The 3-Step Test)} $\rightarrow$ Problems 7, 8, 9, 10, 11, 12
    \item \textbf{Concept 3: Span Identification (Parameter Factoring)} $\rightarrow$ Problems 13, 14, 15, 16, 17
    \item \textbf{Concept 4: Span Membership via Linear Systems} $\rightarrow$ Problems 18, 19, 20, 21, 22
    \item \textbf{Concept 5: Theoretical Properties and Synthesis} $\rightarrow$ Problems 23, 24, 25, 26
    \item \textbf{Concept 6: Diagnostic Testing ("Find the Flaw")} $\rightarrow$ Problems 27, 28, 29, 30
\end{itemize}

\end{document}